% =============================================================================
%  Notas del Profesor — Sesión 04: Representaciones Binarias
%  Programación Avanzada — Universidad Panamericana
%  Compilar: pdflatex sesion04_representaciones_binarias_profesor.tex
% =============================================================================
\documentclass[12pt, oneside]{article}
\usepackage[profesor]{progavanzada}

\begin{document}

\portada{Sesión 04}{Representaciones Binarias}{2026}

\tableofcontents
\newpage

% =============================================================================
\section{Objetivos de la sesión}
% =============================================================================

Al finalizar la sesión el alumno será capaz de:

\begin{enumerate}
  \item Calcular el rango de valores representables con $n$ bits, con y sin signo.
  \item Obtener la representación en complemento a 2 de un entero negativo
        y verificarla.
  \item Interpretar una secuencia de bytes como texto ASCII y viceversa.
  \item Distinguir mayúsculas de minúsculas a nivel de bit en ASCII.
  \item Describir la estructura de un número flotante IEEE 754 y explicar
        por qué ocurren errores de redondeo.
\end{enumerate}

\begin{notaprofesor}[Duración estimada]
  50 minutos. Sugerido: 10 min corrección ejercicios sesión 03 +
  15 min complemento a 2 + 10 min ASCII + 10 min IEEE 754 + 5 min cierre.
\end{notaprofesor}

% =============================================================================
\section{Secuencia de clase}
% =============================================================================

\subsection{Apertura — Corrección de ejercicios (10 min)}

Dedica los primeros minutos a cerrar la sesión anterior.
Los ejercicios 5 y 6 son los que más trabajo dan; revísalos en el pizarrón
mostrando el camino vía binario.

\begin{notaprofesor}[Solución del problema de Carina]
  Los dígitos disponibles son $\{0,1,3,4,5,6,8,9\}$ (8 símbolos = base 8).
  Renombrando $0\to0, 1\to1, 3\to2, 4\to3, 5\to4, 6\to5, 8\to6, 9\to7$,
  Carina cuenta en base 8. El milésimo número:
  \[
    1000_{10} = 1750_8 \;\xrightarrow{\text{remap}}\; \mathbf{1960}
  \]
  Verificación: $1\times512 + 7\times64 + 5\times8 + 0 = 512+448+40 = 1000$. \checkmark\\[4pt]
  Si algún alumno llegó a base 8 por sí solo, aprovecha para reconocerlo:
  es la parte más creativa del problema.
\end{notaprofesor}

\subsection{Complemento a 2 (15 min)}

Antes de explicar el método, lanza la pregunta:

\begin{notaprofesor}[Pregunta de arranque]
  \textit{``Si tengo 8 bits y quiero guardar números negativos,
  ¿cómo lo haría? ¿Cuál bit usaría para el signo?''}\\[4pt]
  La respuesta intuitiva del grupo suele ser ``el bit más a la izquierda
  como signo, el resto como magnitud'' (signo-magnitud). Acepta esa
  respuesta y luego muestra el problema: con signo-magnitud existirían
  \textbf{dos ceros} (\texttt{00000000} y \texttt{10000000}), lo que
  complica la aritmética del procesador. El complemento a 2 elimina ese
  problema y simplifica los circuitos de suma.
\end{notaprofesor}

Desarrolla los tres puntos en pizarrón:
\begin{enumerate}
  \item La tabla de patrones de 8 bits con su valor.
        Énfasis en \texttt{10000000}=$-128$, \texttt{11111111}=$-1$.
  \item El algoritmo invertir+1 con un ejemplo ($-6$, $-50$).
  \item El rango general $-2^{n-1}$ a $+2^{n-1}-1$ y por qué es asimétrico.
\end{enumerate}

\begin{notaprofesor}[Demostración intuitiva del rango asimétrico]
  Pregunta: \textit{``Si niego $-128$, ¿qué obtengo?''}
  Invertir \texttt{10000000} = \texttt{01111111}, sumar 1 = \texttt{10000000}.
  Se regresa a $-128$. Este es el único número que no tiene par positivo.
  Los alumnos suelen encontrar esto perturbador — aprovéchalos: es un
  bug real en software que asume simetría de rango.
\end{notaprofesor}

\subsection{ASCII (10 min)}

Conecta directo con lo anterior:

\begin{notaprofesor}[Puente con complemento a 2]
  \textit{``Ya vimos cómo los números enteros viven en la memoria.
  ¿Cómo vive el texto?''} La respuesta es sencilla: también como números.
  Un carácter = un byte = un entero de 8 bits. ASCII es el acuerdo de
  qué número corresponde a qué carácter.
\end{notaprofesor}

Muestra la tabla A–Z en el pizarrón o proyector. Puntos clave:

\begin{enumerate}
  \item Rango mayúsculas: 65 (A) a 90 (Z), es decir \texttt{0x41}–\texttt{0x5A}.
  \item Rango minúsculas: 97 (a) a 122 (z), es decir \texttt{0x61}–\texttt{0x7A}.
  \item Diferencia exacta de 32 = bit 5 activado.
  \item Ejemplo en vivo: escribe tu nombre en hex. El ejemplo de clase fue ``MAR''
        → \texttt{4D 41 52}.
\end{enumerate}

\begin{notaprofesor}[Actividad rápida de 2 minutos]
  Pide a un voluntario que diga sus iniciales. Calcula con el grupo el hex
  de cada letra. Es inmediato con la tabla y hace el concepto muy concreto.
  También funciona el ejercicio inverso: proyecta una secuencia de bytes
  y pregunta qué dice.
\end{notaprofesor}

\subsection{IEEE 754 (10 min)}

\begin{notaprofesor}[Por qué incluir IEEE 754 aquí]
  Los alumnos declararán \texttt{float} y \texttt{double} desde la
  sesión siguiente. Si no saben que los flotantes tienen errores de
  redondeo inherentes, los bugs que aparecerán en comparaciones y en
  acumulación de errores serán completamente opacos para ellos.
  No es necesario que comprendan la codificación completa; basta con que
  entiendan \textbf{por qué} \texttt{0.1 + 0.2 != 0.3}.
\end{notaprofesor}

Secuencia sugerida:
\begin{enumerate}
  \item Pregunta: \textit{``¿Cómo escribirías 0.1 en binario?''}
        Muestra que $0.1_{10} = 0.0\overline{0011}_2$ (periódico).
        Conecta con $\frac{1}{3} = 0.\overline{3}$ en decimal: no es raro,
        depende de la base.
  \item Muestra la estructura: signo (1 bit), exponente (8 bits),
        mantisa (23 bits) para \texttt{float}.
  \item Regla práctica: \textbf{nunca} comparar flotantes con \texttt{==}.
        Usa \texttt{|a-b| < epsilon}.
\end{enumerate}

\begin{notaprofesor}[Demostración en vivo opcional]
  Si tienes el cuaderno Jupyter abierto, ejecuta:
\begin{verbatim}
#include <iostream>
int main(){
    double a = 0.1 + 0.2;
    std::cout << (a == 0.3) << "\n";   // imprime 0
    std::cout << a << "\n";            // 0.3
    std::cout.precision(17);
    std::cout << a << "\n";            // 0.30000000000000004
}
\end{verbatim}
  El resultado de precisión 17 produce el efecto ``wow'' que fija el
  concepto mejor que cualquier explicación.
\end{notaprofesor}

% =============================================================================
\section{Preguntas frecuentes}
% =============================================================================

\begin{notaprofesor}[FAQ]

\textbf{``¿Por qué el rango de \texttt{char} en C++ va de $-128$ a $127$
y no de $-127$ a $127$?''}\\
Porque \texttt{char} es un entero de 8 bits en complemento a 2.
El patrón \texttt{10000000} se asigna a $-128$ (no existe $-0$).
Nota adicional: en C++ \texttt{char} puede ser \texttt{signed} o
\texttt{unsigned} según la implementación; usa \texttt{int8\_t} /
\texttt{uint8\_t} si necesitas garantías.

\vspace{0.6em}
\textbf{``¿Qué pasa con los acentos y la ñ en ASCII?''}\\
No existen en ASCII estándar (7 bits). El ASCII ``extendido'' de 8 bits
tiene varias codificaciones incompatibles (ISO-8859-1, Windows-1252, etc.).
La solución moderna es \textbf{UTF-8}, que representa los primeros 128
caracteres igual que ASCII y usa 2–4 bytes para los demás.

\vspace{0.6em}
\textbf{``¿Cuándo uso \texttt{float} y cuándo \texttt{double}?''}\\
Por defecto usa \texttt{double}: tiene más precisión (52 bits de mantisa
vs.\ 23) y en CPUs de 64 bits no es más lento. \texttt{float} se usa
en gráficos, audio y aprendizaje automático donde el volumen de datos
hace que el ahorro de memoria importe.

\vspace{0.6em}
\textbf{``¿Existe un estándar para enteros negativos alternativo al
complemento a 2?''}\\
Históricamente sí: signo-magnitud y complemento a 1. Desde C++20
el estándar exige oficialmente complemento a 2 para todos los
enteros con signo. En la práctica, todas las CPUs modernas ya lo usaban.

\end{notaprofesor}

% =============================================================================
\section{Material de la sesión}
% =============================================================================

\begin{itemize}
  \item \textbf{Notas manuscritas originales}: \texttt{materialInicial/sesion04/}
  \item \textbf{Notas alumno}: \texttt{notas\_alumno/pdf/sesion04\_representaciones\_binarias.pdf}
\end{itemize}

% =============================================================================
\section{Rúbrica}
% =============================================================================

\begin{center}
\begin{tabular}{lcc}
  \toprule
  \textbf{Criterio} & \textbf{Puntos} & \textbf{Evidencia} \\
  \midrule
  Convertir entero negativo a complemento a 2 (proceso) & 3 & pasos visibles \\
  Convertir complemento a 2 a decimal (leer negativo)   & 2 & resultado correcto \\
  Decodificar 3 caracteres ASCII desde hex               & 2 & tabla usada correctamente \\
  Codificar nombre/iniciales en hex ASCII                & 2 & coincide con tabla \\
  Identificar campo signo/exponente/mantisa en \texttt{float} & 1 & estructura correcta \\
  \bottomrule
  \textbf{Total} & 10 & \\
\end{tabular}
\end{center}

\end{document}
